%% fuer Drucker "bilby":
%% dvips -O 5mm,1mm

%% fuer Drucker "goofy":
%% dvips600 -O 0mm,10mm Magdeburg

\documentclass[portrait]{seminar}
\usepackage{german,colordvi,boldmath,semcolor,semlayer,rotate,epsf,fancybox,mathsym}
%\input{farbunilogo}
%\input{unilogo}
\usepackage{amssymb}
\UseLegacyTextSymbols

\def\dummy#1{}
\overlaystrue
\def\overlay#1{{}}

\slidewidth9.5in
\slideheight6.7in
%\def\slidetopmargin{\iflandscape.2in\else0.8in\fi}
\def\slidetopmargin{\iflandscape.3in\else0.8in\fi}
\def\slidebottommargin{\iflandscape1.0in\else0.4in\fi}
\def\slideleftmargin{\iflandscape.2in\else.5in\fi}
\def\sliderightmargin{\iflandscape1.0in\else.8in\fi}

\newcommand{\cH}{{\cal H}}

\slidesmag{4}
\centerslidesfalse
\slideframewidth0.5pt
%\slideframe{none}

%\portraitonly
%\landscapeonly
%\def\leftsliderotation#1{\rotate[l]{#1}}
%\sliderotation{left}
\rotateheaderstrue

%%%%%%%%%%%%%%%%%%%%%%%%%%%%%%%%%%%%%%%%%%%%%%
%% Makros
%%%%%%%%%%%%%%%%%%%%%%%%%%%%%%%%%%%%%%%%%%%%%%
\def\bra#1{\left<#1\right|}
\def\ket#1{\left|#1\right>}

\def\trace{\mathop{\rm tr}\nolimits}

\newcommand{\heading}[1]{%
  \begin{center}
    \Large
    \bfseries
    \shadowbox{#1}%
  \end{center}
  \vspace{1ex minus 1ex}%
}
\newbox{\logobox}
%\setbox\logobox=\hbox{\raise 0.17\baselineskip\hbox{\epsfbox{logo.eps}}}

\def\logo{\usebox{\logobox}}
\newpagestyle{MH}{\hss\hbox to \hsize{PhysComp96\hss\thepage}\hss}%
{\hss\hbox to \hsize{${\cal IAKS}$\ \logo\ UniKA\hfil}\hss}

\newpagestyle{nohead}{}{\hss\hbox to \hsize{\tiny IAKS
\ \logo\ Universit"at Karlsruhe\hfil Prof.~Dr.~Th.~Beth}\hss}

\newpagestyle{nohead_beth}{}{\hss\hbox to \hsize{\tiny${\cal IAKS}$\ \logo\
UniKA\hfil}\hss}

\newcommand{\tl}{\tilde}
\newcommand{\onemat}{{\mathbf 1}}

\begin{document}

\begin{slide}
{$ $}
\vfill
\begin{center}
{\Large \blue Mutual simulation of Hamiltonian dynamics\\ on interacting
quantum systems}

\vfill Pawel Wocjan, Dominik Janzing, Martin R{\"o}tteler, Thomas Beth

{\small Institut f"ur Algorithmen und Kognitive Systeme}

{\small Universit"at Karlsruhe}
\end{center}
\vfill
\end{slide}

\begin{slide}
{\large {\bf Control-theoretical model of quantum computers}}
\vfill
{\blue Time evolution} defined by the {\blue
time-dependent Schr\"odinger equation}
\[
\frac{d}{dt} U(t) = -i H(t) U(t)\,,\quad U(0)=\onemat
\]
$H(t)$ Hamiltonian 

$U(t)$ is the resulting unitary after time $t$
\vfill
\[
H(t) = H_d + \sum_{i=1}^m a_i(t) H_i
\]
$H_d$ {\blue drift Hamiltonian} internal to the system

$\sum_{i=1}^m a_i(t) H_i$ {\blue control Hamiltonian} externally
changeable
\vfill
{\red Complexity of a unitary $U$}: minimal time required to implement
$U$
\vfill
\end{slide}

\begin{slide}
{\large {\bf Fast control limit}}
\vfill
strength of the control Hamiltonians arbitrarily large $\Longrightarrow$

{\blue control unitaries in control group $K$} arbitrarily fast
\vfill
{\red Complexity of a unitary $U$}: minimal time
\[
U=\exp(-i H t_N) V_N \cdots \exp(-i H t_2) V_2 \exp(-i H t_1) V_1
\]
\vfill
express it as the evolution according to the piecewise constant Hamiltonian
\[
U = U_N \exp(-i U_N^\dagger H U_N t_N) \cdots \exp(-i U_2^\dagger H U_2 t_2)
    \exp(-i U_1^\dagger H U_1 t_1)
\]
\vfill
\end{slide}


\begin{slide}
{\large {\bf Average Hamiltonian Theory}}
\vfill
{\blue Dyson expansion}
\[
U(t)={\cal T} \exp(\int_{\tau=0}^t d\tau H(\tau))=\exp(-i\bar{H}t)
\]
\vfill
{\blue Magnus expansion} $\bar{H}=\bar{H}^{(0)}+\bar{H}^{(1)}+\cdots$
\begin{eqnarray*}
{\red\bar{H}^{(0)}} & {\red =} & {\red \frac{1}{t}\int_0^t d\tau
H(\tau)\,,\quad\quad\quad\quad\ \quad\quad\quad\quad
\|\bar{H}^{(0)}\|\le 1} \\
\bar{H}^{(1)} & = & \frac{-i}{2t}\int_0^t
d\tau' \int_0^t\tau'' [H(\tau'),H(\tau'')]\,,\quad
\|\bar{H}^{(1)}\|\le t/2
\end{eqnarray*}
\vfill
{\red $\Longrightarrow$ $U(t)$ essentially determined by $\bar{H}^{(0)}$ for
small $t$}
\vfill
\end{slide}


\begin{slide}
{\large {\bf Simulating Hamiltonians (first order simulation)}} 

\vfill
{\blue $H$ can simulate $\tilde{H}$ with overhead $1$}, written
$\tilde{H}\le H$, and {\blue $N$ time steps} iff
\[
\tilde{H}=\sum_{j=1}^N p_j U_j^\dagger H U_j\,,
\]
where $p_1,\ldots,p_N$ are positive real numbers summing to $1$ and
$U_j\in K$ are control operations, i.e.
{\red $\tilde{H}$ is a convex sum of unitary conjugates of $H$}
\vfill
$H$ can simulate $\tilde{H}$ with overhead {\red $\tau$} iff 
$\tilde{H} \le {\red \tau} H$ 
\vfill 
{\red time-optimal simulation} problem $\longrightarrow$
{\red convex optimization} problem 
\vfill
\end{slide}

\begin{slide}
{\large {\bf Composite systems}}
\vfill
system consisting of {\blue $n$ qudits}  
$\quad {\cal H}=\C^d\otimes\C^d\otimes\cdots\otimes\C^d$ 
\vfill
{\blue pair-interaction Hamiltonian}
\[
H = 
\sum_{k<l} \sum_{\alpha\beta} J_{kl;\alpha\beta} 
\sigma_\alpha^{(k)}\sigma_\beta^{(l)} +
\sum_k \sum_{\alpha} r_{k;\alpha} \sigma_\alpha^{(k)}\,,
\]
where $B:=\{\sigma_\alpha\mid\alpha=1,\ldots,d^2-1\}$ is a basis of $su(d)$

%{\blue coupling matrix}
%\[
%J=\left(
%\begin{array}{c|c|c|c|c}
%0      & J_{12} & J_{13} & \cdots & J_{1n} \\ \hline
%J_{21} & 0      & J_{23} & \cdots & J_{2n} \\ \hline
%\vdots & \vdots &        &        &        \\ \hline
%J_{n1} & J_{n2} & J_{n3} &        & 0 
%\end{array}
%\right)\in\R^{(d^2-1)\,n\times (d^2-1)\,n}
%\]

\vfill
{\blue control operations} are all {\red local} operations in
\[
K:=SU(d)\otimes SU(d)\otimes\cdots\otimes SU(d)
\]
\vfill
\end{slide}

\begin{slide}
{\blue Decoupling/Selective decoupling}

\quad quantum computing requires the ability to perform gates

\quad couplings created by gates can only originate from natural couplings

\quad naturally available Hamiltonians do not couple specific pairs of qubits

\quad {\red $\longrightarrow$ turn off disturbing couplings}

\vfill
{\blue Time-reversal}

\quad evolution of a single spin seems is phenomenologically a relaxation/

\quad dephasing process

%\quad (described by Bloch equations)

\quad however, it is not always irreversible

\quad {\red $\longrightarrow$ refocusing}\quad simulate $-H$ when $H$ is
the system Hamiltonian

\vfill 
{\blue Universal simulation of quantum dynamics}

\quad simulating quantum dynamics on a classical computer is hard

\quad {\red $\longrightarrow$ universal quantum simulator}

\quad\quad\quad use a quantum computer to simulate arbitrary dynamics
\vfill
\end{slide}


\begin{slide}
{\bf Algebraic Graph Theory gives Lower Bounds on Complexity}
\vfill
represent {\blue pair-interaction Hamiltonian}
\[
H = 
\sum_{k<l} \sum_{\alpha\beta} J_{kl;\alpha\beta} 
\sigma_\alpha^{(k)}\sigma_\beta^{(l)} +
\sum_k \sum_{\alpha} r_{k;\alpha} \sigma_\alpha^{(k)}\,,
\]
where $B:=\{\sigma_\alpha\mid\alpha=1,\ldots,d^2-1\}$ is a ONB of $su(d)$
\vfill
as {\blue coupling matrix} (extension of adjacency matrix of a graph)
\[
J=\left(
\begin{array}{c|c|c|c|c}
0      & J_{12} & J_{13} & \cdots & J_{1n} \\ \hline
J_{21} & 0      & J_{23} & \cdots & J_{2n} \\ \hline
\vdots & \vdots &        &        &        \\ \hline
J_{n1} & J_{n2} & J_{n3} &        & 0 
\end{array}
\right)\in\R^{(d^2-1)\,n\times (d^2-1)\,n}
\]
\vfill
\end{slide}

\begin{slide}
{\bf Analogy: coupling matrix $\leftrightarrow$ adjacency matrix}
\vfill
consider dipole-dipole coupling between all nodes
\[
H=\sum_{k<l} w_{kl} (3\,\sigma_z\otimes\sigma_z - \sum_{\alpha=x,y,z} 
\sigma_\alpha\otimes\sigma_\alpha)\,,\quad w_{kl}\neq 0
\]
\vfill
coupling matrix can be expressed as
\[
J=W\otimes C\,,
\]
where $W=(w_{kl})$ and $C={\rm diag}(-1,-1,2)$.
\vfill
$W$ adjacency matrix of a weighted graph

$C$ characterizes the type of coupling between nodes
\vfill
\end{slide}


\begin{slide}
{\large {\bf Action of control operations on the coupling matrix $J$}}
\vfill
\begin{center}
{\blue Conjugation of the pair-interaction Hamiltonian $H$ by control
operations}

corresponds to

{\blue conjugation of $J$ by a direct sum of orthogonal matrices}
\end{center}
\[
\begin{array}{ccc}
H & \mapsto & (U_1\otimes U_2\otimes\cdots\otimes U_n)^\dagger\, H\, 
(U_1\otimes U_2\otimes\cdots\otimes U_n) \\
J & \mapsto & (O_1\oplus O_2\oplus\cdots\oplus O_n)\, J\,
(O_1\oplus O_2\oplus\cdots\oplus O_n)^T
\end{array}
\]
\vfill
\end{slide}

\begin{slide}
{\large {\bf Lower bounds on time overhead}}
\vfill
If $H$ can simulate $\tilde{H}$ with overhead $\tau$ then
\[
\tilde{J}/\tau = \sum_{j=1}^N p_j 
(U_{j1}\oplus U_{j2}\oplus\cdots\oplus U_{jn})
\,J\, (U_{j1}\oplus U_{j2}\oplus\cdots\oplus U_{jn})^T
\]
convex sum of conjugates of $J$
\vfill
$\Rightarrow$ necessary condition given by {\blue majorization criterion}:\\
\begin{center}
the spectrum Spec($\tilde{J}$) majorizes the spectrum Spec($\tau J$)
\end{center}
\vfill
\end{slide}

\begin{slide}
{\large {\bf Application I: Time reversal}}
\vfill
Time reversal: $H \mapsto -H$ (refocusing, spin echoes)
\vfill
dipole-dipole coupling between all nodes (easy to invert)
\begin{eqnarray*}
H & = & \sum_{k<l} (3\,\sigma_z\otimes\sigma_z - \sum_{\alpha=x,y,z} 
\sigma_\alpha\otimes\sigma_\alpha) \\
J & = & K \otimes {\rm diag}(-1,-1,2)\,,
\end{eqnarray*}
where $K$ is the adjacency matrix of the complete graph
\vfill
lower bound on time overhead is $2$
\vfill
this bound is tight (well known-scheme in NMR)
\vfill
\end{slide}


\begin{slide}
{\large {\bf Application I: Time reversal}}
\vfill
strong scalar coupling between all nodes (difficult to invert)
\begin{eqnarray*}
H & = & \sum_{k<l}
(\sigma_x^{(k)}\sigma_x^{(l)} + \sigma_y^{(k)}\sigma_y^{(l)} + 
 \sigma_z^{(k)}\sigma_z^{(l)})\\
J & = & K \otimes {\bf 1}_3
\end{eqnarray*}
\vfill
lower bound on the overhead $n-1$

the inverted time-evolution is necessarily slower by the factor $(n-1)$
\vfill
upper bound $cn-1$ (decoupling scheme based on orthogonal arrays)
\vfill
\end{slide}

\begin{slide}
{\bf Decoupling and time-reversal schemes (Upper bounds)}
\vfill 
Conjugating by unitaries of an {\blue error basis} ${\cal E}=\{U_1={\mathbf
1},U_2,\ldots,U_{d^2}\}$
\[
\sum_j U_j^\dagger\, a\, U_j = {\mathbf 0}\quad\quad
\mbox{for all } a\in su(d)
\]
cancels all matrices in $su(d)$
\vfill
Combinatorial problem: which sequences of conjugations allow to switch off
all interactions?
\vfill
Efficient solution: decoupling in networks based on
{\blue orthogonal arrays} 

number of conjugations grows linearly with $n$
\vfill
\end{slide}

\begin{slide}
{\large {\bf Application II: Hierarchy of Hamiltonians}}
\vfill
\begin{eqnarray*}
H_1 & = & \sum_{k<l} w_{kl} 
(\sigma^{(k)}_x\otimes\sigma^{(l)}_x\, {\red +}\, 
 \sigma^{(k)}_y\otimes\sigma^{(l)}_y) \\
& & \\
H_2 & = & \sum_{k<l} w_{kl} 
(\sigma^{(k)}_x\otimes\sigma^{(l)}_x\, {\blue -}\, 
 \sigma^{(k)}_y\otimes\sigma^{(l)}_y)
\end{eqnarray*}
\vfill
$H_1$ cannot simulate $H_2$ with overhead less than $(n-1)$

$H_2$ can simulate $H_1$ with overhead $2$
\vfill
\end{slide}

\begin{slide}
{\large {\bf Application III: Switching off {\red some} interactions}}
\vfill
Combinatorial problem: cancel some interactions
\vfill
Are the other interactions necessarily disturbed?
\begin{figure}
\centerline{
\epsfbox[0 0 285 98]{HonnefGraphen.eps}}
\end{figure}
\vfill 
Yes, for some graphs $\rightarrow$ lower bound on time overhead
given by eigenvalues of the adjacency matrix \vfill

Upper bound given by chromatic index of interaction graph
\end{slide}

\begin{slide}
{\large {\bf Universal simulation using finite control groups}}
\vfill
{\blue Universal simulation} does not require all operations in $SU(d)$

Certain subgroups are sufficient; called {\blue transformer groups}
\vfill 
{\blue Finite transformer groups are characterized in terms of group 
characters} 
\vfill
\end{slide}

\begin{slide}
{\bf Treating the Independent Set Problem by $2$D Ising interactions with
Adiabatic Quantum Computing}
\vfill
construct a {\red ``physically reasonable''} Hamiltonian whose ground states
encode the solutions of Independent Set:
\bigskip
\begin{itemize}
\item INSTANCE: Graph $G=(V,E)$, positive integer $v\le |V|$.
\item QUESTION: Does $G$ contain an independent set whose cardinality
is at least $v$, i.e., a subset $V'\subseteq V$ such that $|V'|\ge v$
and such that no two vertices in $V'$ are joined by an edge in $E$?
\end{itemize}
Independent Set remains NP-complete for {\blue cubic planar} graphs.
\vfill
\end{slide}

\begin{slide}
Cubic planar graph
\centerline{
\epsfbox[20 25 130 130]{graph3.ps}}
{\bf Theorem (Planar spin glass within a magnetic field)}\\
Let $G=(V,E)$ be a cubic planar. Determining the energy of the ground
states of the corresponding Hamiltonian
\begin{equation}
H_P=\sum_{(k,l)\in E} \sigma_z^{(k)}\sigma_z^{(l)} + \sum_{k\in V}
\sigma_z^{(k)}
\end{equation}
is equivalent to determining the maximum cardinality of independent
sets of $G$.

Problem: $H_P$ is {\red not physically reasonable}
\end{slide}

\begin{slide}
{\bf Planar orthogonal embedding} of a graph $G=(V,E)$ maps
\begin{itemize}
\item vertices $k\in V$ to lattice points $\Gamma(k)$ and 
\item edges $(k,l)\in E$ to paths in the lattice such that the images of
their endpoints $\Gamma(k)$ and $\Gamma(l)$ are connected and such
that the paths do not share any vertices (besides the endpoints).
\end{itemize}
\centerline{
\epsfbox[80 80 160 160]{Grafik2.eps}}
\end{slide}

\begin{slide}
{\bf Planar orthogonal Hamiltonians}

a physically reasonable Hamiltonian whose ground states
correspond $1$-to-$1$ to the grounds states of $H_P$

\centerline{
\epsfbox[80 80 100 160]{Grafik1.eps}}
\bigskip

\bigskip

\bigskip

\bigskip
filled circles: dummies, no local $\sigma_z$ terms

thick lines: ${\red -c}\sigma_z\otimes\sigma_z$ coupling

circles: vertices $\Gamma(k)$, $\sigma_z$ as local term

thin lines: $\sigma_z\otimes\sigma_z$ coupling
\end{slide}



\begin{slide}
{\bf Mathematical tools}

{\blue Vector majorization}

let $x=(x_1,x_2,\ldots,x_d)\in\R^d$

define
$x^\downarrow=(x^\downarrow_1,x^\downarrow_2,\ldots,x^\downarrow_d)\in\R^d$
as the vector whose components are the same as of $x$, but ordered so that
\[
x^\downarrow_1\ge x^\downarrow_2\ge\ldots\ge x^\downarrow_d
\]

{\blue $x$ majorizes $y$}, written {\blue $x\preceq y$}, if
\[
\sum_{j=1}^k x^\downarrow_j \le \sum_{j=1}^k y^\downarrow_j
\]
for $k=1,\ldots,d-1$ and 
\[
\sum_{j=1}^d x^\downarrow_j = \sum_{j=1}^d y^\downarrow_j
\]
\end{slide}

\begin{slide}
{\bf Mathematical tools} 

{\blue Operator majorization}
\vfill
let $A$ and $B$ be two Hermitian matrices of size $d$

we say {\blue $A$ majorizes $B$}, written {\blue $A\preceq B$}, if
\[
\lambda(A)\preceq\lambda(B)\,,
\]
where $\lambda(A)$ and $\lambda(B)$ are the spectra (the vectors of
eigenvalues) of $A$ and $B$, respectively.
\vfill
{\bf Uhlmann's theorem} 

we have $A\preceq B$ iff
\[
A=\sum_{j=1}^N p_j U_j^\dagger B U_j\,,
\]
where $\{p_j\}$ is a probability distribution and $\{U_j\}$ is a set
of unitaries
\vfill
\end{slide}

\end{document}









